% Cover letter for the ACM TOPS submission. Standalone: compiles on its own with
%   tectonic -X compile cover_letter.tex
% and does not depend on main.tex or references.bib. Works cited live in cover_letter.bib.
\documentclass[11pt,a4paper]{article}

\usepackage[T1]{fontenc}
\usepackage[utf8]{inputenc}
\usepackage{libertine}
\usepackage[margin=1in]{geometry}
\usepackage{microtype}
\usepackage[hidelinks]{hyperref}
\usepackage{parskip}
\usepackage[numbers,sort&compress]{natbib}

\setlength{\bibsep}{2pt plus 0.3ex}
\renewcommand{\bibsection}{\subsection*{\normalsize References}}
\pagestyle{empty}

% Run-in headings keep this reading as a letter rather than a paper.
\newcommand{\lhead}[1]{\medskip\noindent\textbf{#1}\quad\ignorespaces}

\begin{document}

\noindent
Professor Michael Waidner\\
Editor-in-Chief, \emph{ACM Transactions on Privacy and Security}\\
ATHENE $|$ TU Darmstadt $|$ Fraunhofer SIT

\medskip
\noindent\today

\bigskip
\noindent Dear Professor Waidner,

\medskip

We submit \emph{Topology-Aware Differential Privacy in Hierarchical Federated Learning} for
consideration as a regular paper in \emph{ACM Transactions on Privacy and Security}.

\lhead{The problem.}
Hierarchical federated learning places regional aggregators between clients and the cloud, so
that a participant's update reaches any observer above that tier only after it has been
combined with those of its neighbours. The concealment this provides is not uniform. It
depends on the effective size of the aggregation region, and regions in operational
deployments differ by an order of magnitude. Prevailing practice nevertheless applies a single
noise multiplier to every participant, calibrated for the most exposed region. Every other
participant then carries more noise than its own exposure requires, and the federation pays
for that surplus in accuracy.

\lhead{What we contribute.}
We give a silo-level differential privacy guarantee for the mechanism, and separately bound
the mutual information between a participant's local class distribution and any estimate an
observer above the regional tier can form. The two results are kept apart deliberately: they
rest on different assumptions and answer different questions, and neither is derived from the
other. Throughout, the adjacency used by the mechanism is matched to the granularity of the
property being protected, which is a property of a participant's dataset as a whole rather
than of any individual record.

Minimising the worst-case value of that bound under a fixed utility budget yields an explicit,
min-max optimal allocation. The budget it recovers relative to uniform noise has a closed
form, $\delta = 1 - \langle \rho \rangle_V / \rho_{\max}$, which depends only on the region
structure and the aggregation weights. A practitioner can therefore determine before training
whether the method will help, and $\delta$ is exactly zero when all regions are equally
exposed, which supplies an explicit applicability test in place of a claim of universal
benefit. Experiments on image and text classification confirm the prediction at
$\varepsilon = 0.99$, including a balanced control deployment for which the theory predicts
parity with the baseline and which delivers it.

\lhead{Why this belongs in TOPS.}
Two lines of work in the journal bear directly on this paper, and to our reading they have so
far run in parallel.

The first is differentially private mechanism design with formal guarantees, represented by
Zhang et al.\ on real-time release of sequential data \cite{zhang2023sequential} and by
Clifton et al.\ on private $k$-nearest-neighbour imputation \cite{clifton2022knn}. That work
treats calibration rigorously, but outside federated settings.

The second is privacy and security in federated learning, which has grown considerably in the
journal: poisoning conducted under the cover of differential privacy
\cite{zheng2025dppoison}, backdoor attacks in peer-to-peer federated learning
\cite{syros2025backdoor}, defences against poisoning \cite{zheng2025gradcam}, and vertical
federated learning for private data publication \cite{xun2026vflgan}. That work is
predominantly adversarial, and treats the calibration of the noise as given.

Our paper joins the two. It asks how differentially private noise should be calibrated to the
structure of a federated deployment, and answers in closed form. The relationship to
Zheng et al.\ \cite{zheng2025dppoison} is the most direct. DP-Poison shows that an adversary
can conceal a poisoning attack within the noise that differential privacy introduces, treating
the magnitude of that noise as a fixed quantity. We ask the prior question of how much noise
each participant should carry, and show that the prevailing answer systematically
over-protects most of the federation.

\lhead{Originality and prior dissemination.}
The manuscript is original, is not under review at any other venue, and has not been published
in whole or in part.
% ACTION FOR THE AUTHORS: ACM's prior-publication policy requires disclosing any publicly
% posted preprint. If an earlier version of this work is posted, state it here with its
% identifier and delete this comment. If none is posted, delete this comment only.

\lhead{Reproducibility.}
Every experimental artefact is archived, and each table and figure in the paper is regenerated
by a script in the accompanying repository. A verification script re-derives every reported
privacy parameter from the mechanism itself rather than reading a stored value, so the
guarantees we report can be checked independently of our own tooling.

\medskip
We look forward to the reviewers' assessment.

\bigskip
\noindent Sincerely,

\medskip
\noindent
Murtaza Rangwala, Richard O. Sinnott, and Rajkumar Buyya\\
School of Computing and Information Systems\\
The University of Melbourne, Australia

\bigskip
\bibliographystyle{unsrtnat}
\bibliography{cover_letter}

\end{document}
